% ============================================================
%  Documentation — Générateur de molécules non-benzenoïdes
%  Stage AMU — Master 2 IMD / CSP
% ============================================================
\documentclass[a4paper,12pt]{article}

\usepackage[utf8]{inputenc}
\usepackage[T1]{fontenc}
\usepackage[french]{babel}
\usepackage{geometry}
\geometry{margin=2.5cm}

\usepackage{amsmath, amssymb}
\usepackage{graphicx}
\usepackage{hyperref}
\usepackage{xcolor}
\usepackage{listings}
\usepackage{booktabs}
\usepackage{array}
\usepackage{enumitem}
\usepackage{titlesec}
\usepackage{fancyhdr}
\usepackage{tcolorbox}
\usepackage{float}
\usepackage{microtype}
\usepackage{parskip}

% ---- couleurs ----
\definecolor{codebg}{rgb}{0.96,0.96,0.96}
\definecolor{codeframe}{rgb}{0.75,0.75,0.75}
\definecolor{primary}{RGB}{0,90,160}
\definecolor{accent}{RGB}{200,60,0}
\definecolor{notegreen}{RGB}{0,120,60}

% ---- listings (Python) ----
\lstdefinestyle{python}{
  language=Python,
  backgroundcolor=\color{codebg},
  frame=single,
  rulecolor=\color{codeframe},
  basicstyle=\ttfamily\small,
  keywordstyle=\bfseries\color{primary},
  commentstyle=\itshape\color{notegreen},
  stringstyle=\color{accent},
  numbers=left,
  numberstyle=\tiny\color{gray},
  numbersep=8pt,
  breaklines=true,
  showstringspaces=false,
  tabsize=4,
}
\lstset{style=python}

% ---- boîtes ----
\tcbuselibrary{skins, breakable}
\newtcolorbox{remarque}[1][]{
  colback=blue!5, colframe=primary, fonttitle=\bfseries,
  title=Remarque, #1
}
\newtcolorbox{attention}[1][]{
  colback=orange!8, colframe=accent, fonttitle=\bfseries,
  title=Attention, #1
}
\newtcolorbox{todo}[1][]{
  colback=gray!8, colframe=gray!60, fonttitle=\bfseries,
  title=À faire, #1
}

% ---- en-têtes ----
\pagestyle{fancy}
\fancyhf{}
\rhead{\textcolor{gray}{\small Générateur non-benzenoïde — Stage AMU M2}}
\lhead{\textcolor{gray}{\small \leftmark}}
\cfoot{\thepage}

\hypersetup{
  colorlinks=true,
  linkcolor=primary,
  urlcolor=primary,
  citecolor=primary
}

% ============================================================
\begin{document}

\begin{titlepage}
  \centering
  \vspace*{2cm}
  {\huge\bfseries\color{primary}
    Générateur de Structures Moléculaires\\[0.4em]
    Non-Benzenoïdes Polycycliques\\[1em]}
  {\large\color{gray} Documentation technique — Stage AMU, Master 2 IMD}\\[0.5em]
  {\large\color{gray} Projet CSP — Table de voisinage}\\[3em]
  \rule{\linewidth}{0.5pt}\\[1em]
  {\large
    \textbf{Contexte :} génération automatique de molécules polycycliques\\
    pour alimenter un modèle de satisfaction de contraintes (Choco Solver)
  }\\[1em]
  \rule{\linewidth}{0.5pt}\\[3em]
  {\normalsize
    \textbf{Implémentation :} Python 3 (bibliothèques standard uniquement)\\[0.5em]
    \textbf{Format de sortie :} Chemical Markup Language (CML)\\[0.5em]
    \textbf{Date :} Avril 2026
  }
  \vfill
\end{titlepage}

\tableofcontents
\newpage

% ============================================================
\section{Contexte et objectif}
% ============================================================

\subsection{Motivation}

Ce projet s'inscrit dans la conception d'un solveur de contraintes (Choco Solver) capable de
générer des molécules polycycliques planaires. Pour alimenter ce solveur, il est nécessaire de
disposer d'une \textbf{table de voisinage} exhaustive : pour chaque type de cycle (pentagone,
hexagone, heptagone), quels sont tous les arrangements valides de cycles voisins ?

Le présent générateur a pour rôle de produire automatiquement, en Python, \textbf{toutes les
structures moléculaires distinctes} de la famille des hydrocarbures polycycliques non-benzenoïdes,
et de les exporter au format CML pour utilisation ultérieure.

\subsection{Famille de molécules ciblée}

On s'intéresse aux \textbf{hydrocarbures polycycliques} (composés uniquement de C et H) dont les
cycles sont de taille 5, 6 ou 7. Ces molécules sont dites \textit{non-benzenoïdes} car elles
contiennent au moins un cycle différent d'un hexagone (à la différence du benzène et de ses
dérivés purement hexagonaux).

Chaque molécule est composée :
\begin{itemize}
  \item d'un \textbf{cycle central} (pentagone, hexagone ou heptagone) ;
  \item d'un ou plusieurs \textbf{cycles voisins}, fusionnés sur les arêtes du cycle central.
\end{itemize}

Les cycles sont fusionnés par \textbf{partage d'arête} (\textit{edge-fused}), ce qui est la
topologie standard des HAP (hydrocarbures aromatiques polycycliques). Chaque carbone respecte une
valence totale de 4.

% ============================================================
\section{Contraintes de génération}
% ============================================================

\subsection{Représentation par séquence de voisins}

Une molécule est entièrement codée par le couple $(\text{taille\_centrale},\ s)$ où :
\[
  s = (s_0,\ s_1,\ \ldots,\ s_{n-1})\ \in \{0, 5, 6, 7\}^n
\]
$n$ est la taille du cycle central, et $s_i$ est la taille du cycle voisin fusionné sur l'arête
$i$ du cycle central ($s_i = 0$ signifie qu'aucun voisin n'est présent sur cette arête).

\begin{remarque}
  Les arêtes du cycle central sont numérotées de 0 à $n-1$ dans l'ordre trigonométrique
  (sens anti-horaire). L'arête $i$ relie le sommet $i$ au sommet $(i+1) \bmod n$.
\end{remarque}

\subsection{Contraintes sur le nombre de voisins}

Toutes les positions ne peuvent pas être remplies pour toutes les tailles de cycle central.
Les limites sont les suivantes :

\begin{center}
\begin{tabular}{cccp{7cm}}
  \toprule
  Cycle central & Min voisins & Max voisins & Raison du maximum \\
  \midrule
  5 (pentagone)  & 2 & 4 & 5 voisins fermeraient un hexagone intérieur complet — géométriquement impossible sans distorsion majeure \\
  6 (hexagone)   & 2 & 6 & toutes les arêtes peuvent être fusionnées \\
  7 (heptagone)  & 2 & 5 & au-delà de 5 voisins, les contraintes angulaires créent des chevauchements géométriques \\
  \bottomrule
\end{tabular}
\end{center}

On exclut également les séquences \textbf{purement benzenoïdes} : un hexagone central entouré
uniquement de cycles à 6 appartient à la famille du coronène et de ses homologues, déjà couverts
par la chimie classique.

\subsection{Symétries : le groupe diédral $D_n$}

Deux séquences $s$ et $s'$ décrivent la \textbf{même molécule} si l'une peut être obtenue à
partir de l'autre par une rotation ou une réflexion du cycle central. Ces transformations
constituent le \textbf{groupe diédral} $D_n$, qui contient $2n$ éléments.

\subsubsection{Rotations}

Une rotation d'angle $k \times \frac{2\pi}{n}$ ($k = 0, 1, \ldots, n-1$) correspond sur la
séquence à un \textbf{décalage circulaire} :
\[
  R_k(s)_i = s_{(i+k) \bmod n}
\]

Par exemple, pour $n = 5$ :
\[
  (5, 6, 5, 0, 0) \xrightarrow{R_1} (6, 5, 0, 0, 5) \xrightarrow{R_2} (5, 0, 0, 5, 6) \xrightarrow{\cdots}
\]
Ces quatre séquences représentent la même molécule vue sous différentes orientations.

\subsubsection{Réflexions (miroirs)}

Une réflexion selon un axe de symétrie du polygone correspond sur la séquence à une
\textbf{inversion de l'ordre de parcours}. Les $n$ réflexions s'écrivent :
\[
  M_k(s)_i = s_{(k - i) \bmod n}
\]

Une molécule et son image miroir sont donc considérées comme identiques (on ne distingue pas
les énantiomères ici — les structures planes n'ont pas de chiralité intrinsèque).

\subsubsection{Taille du groupe}

\begin{center}
\begin{tabular}{ccc}
  \toprule
  Cycle central & Taille de $D_n$ & Nombre de transformations \\
  \midrule
  5 & $|D_5| = 10$ & 5 rotations + 5 réflexions \\
  6 & $|D_6| = 12$ & 6 rotations + 6 réflexions \\
  7 & $|D_7| = 14$ & 7 rotations + 7 réflexions \\
  \bottomrule
\end{tabular}
\end{center}

\subsection{Forme canonique}

Pour éviter les doublons, on ne conserve qu'un seul représentant par orbite sous $D_n$. On
choisit la \textbf{forme canonique} définie comme suit :

\begin{enumerate}
  \item Parmi toutes les $2n$ transformées de $s$, ne retenir que celles dont les zéros sont
    \textbf{groupés à droite} (toutes les positions vides sont consécutives en fin de séquence).
    Cette convention place les voisins en premier, dans un ordre déterministe.
  \item Parmi ces candidats, choisir le \textbf{maximum lexicographique} (les grandes valeurs
    en tête).
\end{enumerate}

\begin{remarque}
  La contrainte ``zéros à droite'' garantit que le bloc de voisins non nuls est toujours
  aligné à gauche, ce qui facilite la comparaison et l'indexation dans la table CSP.
  Le maximum lexicographique brise les symétries résiduelles de manière unique.
\end{remarque}

Si aucune transformée ne satisfait la condition des zéros regroupés (cas des séquences
entièrement non nulles, ex. hexagone complet), on prend simplement le maximum lexicographique
de toutes les transformées.

\subsection{Dénombrement exact : lemme de Burnside}

Le nombre de séquences canoniques (orbites distinctes sous l'action de $D_n$) se calcule
exactement grâce au \textbf{lemme de Burnside} :
\[
  |\text{orbites}| = \frac{1}{|D_n|} \sum_{g \in D_n} |\text{Fix}(g)|
\]
où $\text{Fix}(g)$ désigne l'ensemble des séquences valides qui sont \textit{invariantes} par
la transformation $g$.

\begin{remarque}[title=Intuition du lemme de Burnside]
  Le lemme de Burnside (aussi appelé lemme de Cauchy--Frobenius) affirme que le nombre d'objets
  \textbf{distincts sous un groupe de symétries} $G$ est la \textbf{moyenne} du nombre de points
  fixes de chaque symétrie. Autrement dit : pour compter combien de ``colliers'' différents on
  peut fabriquer avec des perles colorées, on regarde pour chaque rotation et chaque réflexion
  combien de colliers restent inchangés, puis on fait la moyenne. Ici, les ``colliers'' sont les
  séquences de voisins $(s_0, \ldots, s_{n-1})$ et le groupe est $D_n$.
\end{remarque}

\subsubsection{Séquences brutes valides}

L'ensemble de départ $S$ contient les séquences $(s_0, \ldots, s_{n-1}) \in \{0,5,6,7\}^n$
vérifiant :
\begin{enumerate}
  \item le nombre de positions non nulles est compris entre $k_{\min}$ et $k_{\max}$ ;
  \item \textbf{exclusion benzenoïde} (cycle central 6 uniquement) : au moins un voisin
    non nul doit être $\neq 6$.
\end{enumerate}

Le cardinal brut (avant réduction par symétrie) est :
\[
  |S| = \sum_{k=k_{\min}}^{k_{\max}} \binom{n}{k} \cdot 3^k
  \quad - \quad \underbrace{\sum_{k=k_{\min}}^{k_{\max}} \binom{n}{k}}_{\text{benzenoïdes (si $n=6$)}}
\]
Le facteur $3^k$ correspond aux 3 choix $\{5, 6, 7\}$ pour chaque position non nulle.

\subsubsection{Points fixes par type de transformation}

\paragraph{Cas $n$ premier (5 et 7).}
Comme $n$ est premier, toute rotation non triviale $r^k$ ($k \neq 0$) a pour seuls points
fixes les séquences \textbf{constantes} $(v, v, \ldots, v)$ avec $v \in \{0,5,6,7\}$.
Or une séquence constante non nulle a $n$ positions occupées, ce qui dépasse le maximum autorisé
($k_{\max} = 4$ pour $n=5$, $k_{\max} = 5$ pour $n=7$). Donc :
\[
  \text{Fix}(r^k) = 0 \quad \text{pour } k = 1, \ldots, n-1
\]

Chaque réflexion fixe 1 position et échange $\frac{n-1}{2}$ paires. Par exemple, pour $n=5$,
la réflexion $\sigma_0$ fixe la position 0 et échange $(1 \leftrightarrow 4)$ et
$(2 \leftrightarrow 3)$. Les paramètres libres sont $(s_0, s_1, s_2)$, et le comptage des
configurations valides donne $\text{Fix}(\sigma) = 33$ (pour $n=5$) et $\text{Fix}(\sigma) = 144$
(pour $n=7$).

\paragraph{Cas $n = 6$ (non premier).}
Les rotations d'ordre strict inférieur à 6 ont des points fixes non triviaux :

\begin{center}
\begin{tabular}{llcl}
  \toprule
  Transformation & Orbites sur les positions & Ordre & $|\text{Fix}|$ \\
  \midrule
  $r, r^5$   & $\{0,1,2,3,4,5\}$ & 6 & 2 (séquences constantes $(5,\ldots)$ et $(7,\ldots)$) \\
  $r^2, r^4$ & $\{0,2,4\},\ \{1,3,5\}$ & 3 & 12 \\
  $r^3$      & $\{0,3\},\ \{1,4\},\ \{2,5\}$ & 2 & 56 \\
  \bottomrule
\end{tabular}
\end{center}

De plus, $D_6$ possède \textbf{deux classes de conjugaison} de réflexions :
\begin{itemize}
  \item 3 réflexions par sommets opposés (fixent 2 positions, échangent 2 paires) :
    $|\text{Fix}| = 236$ chacune ;
  \item 3 réflexions par milieux d'arêtes opposées (échangent 3 paires) :
    $|\text{Fix}| = 56$ chacune.
\end{itemize}

\subsubsection{Résultats}

\paragraph{Cycle 5 ($|D_5| = 10$, $|S| = 765$) :}
\[
  |\text{orbites}| = \frac{1}{10}\left(765 + 4 \times 0 + 5 \times 33\right)
                   = \frac{930}{10} = \mathbf{93}
\]

\paragraph{Cycle 6 ($|D_6| = 12$, $|S| = 4020$ après exclusion benzenoïde) :}
\[
  |\text{orbites}| = \frac{1}{12}\left(
    4020 + 2 + 12 + 56 + 12 + 2 + 3 \times 236 + 3 \times 56
  \right)
  = \frac{4980}{12} = \mathbf{415}
\]

\paragraph{Cycle 7 ($|D_7| = 14$, $|S| = 9072$) :}
\[
  |\text{orbites}| = \frac{1}{14}\left(9072 + 6 \times 0 + 7 \times 144\right)
                   = \frac{10080}{14} = \mathbf{720}
\]

\begin{center}
\begin{tabular}{ccccc}
  \toprule
  Cycle central & $|S|$ brut & $|D_n|$ & Orbites & Voisins (min--max) \\
  \midrule
  5 &  765 & 10 &  93 & 2--4 \\
  6 & 4020 & 12 & 415 & 2--6 \\
  7 & 9072 & 14 & 720 & 2--5 \\
  \midrule
  \textbf{Total} & & & \textbf{1228} & \\
  \bottomrule
\end{tabular}
\end{center}

\begin{remarque}
  Le fait que chaque résultat tombe sur un entier exact constitue une vérification indirecte
  de la cohérence entre l'implémentation du code (énumération + forme canonique) et la théorie
  (lemme de Burnside). Toute erreur dans les contraintes ou dans le groupe de symétrie
  produirait un résultat fractionnaire.
\end{remarque}

% ============================================================
\section{Concepts géométriques clés}
% ============================================================

\subsection{Polygone régulier et rayons caractéristiques}

Tous les cycles sont traités comme des \textbf{polygones réguliers} de côté $d = 1.42\ \text{\AA}$
(distance C--C standard dans les HAP aromatiques). Deux rayons sont utilisés :

\begin{itemize}
  \item \textbf{Rayon circonscrit} (sommet $\to$ centre) :
    \[
      R_n = \frac{d}{2 \sin(\pi/n)}
    \]
  \item \textbf{Apothème} (centre $\to$ milieu d'arête) :
    \[
      a_n = \frac{d}{2 \tan(\pi/n)}
    \]
\end{itemize}

Le centre d'un cycle voisin fusionné sur l'arête $[v_1, v_2]$ est placé à distance $a_n$ du
milieu de cette arête, du côté \textit{extérieur} au cycle central.

\subsection{Sens de placement des sommets}

À partir du centre $C$ du nouveau cycle et de $v_1$ (premier sommet de l'arête partagée),
les sommets sont placés en \textbf{sens horaire} :
\[
  P_k = C + R_n \cdot \left(\cos\left(\theta_0 - k \cdot \frac{2\pi}{n}\right),\
                             \sin\left(\theta_0 - k \cdot \frac{2\pi}{n}\right)\right)
  \quad k = 0, 1, \ldots, n-1
\]
où $\theta_0 = \text{atan2}(v_{1y} - C_y,\ v_{1x} - C_x)$. Le signe négatif ($-k$) assure
le sens horaire, qui place $v_2$ en position $k=1$, cohérent avec l'orientation de l'arête
$v_1 \to v_2$.

\subsection{Fusion coadjacente}

Quand deux voisins occupent des positions \textbf{consécutives} $i$ et $i+1$, leurs cycles ne
partagent pas seulement une arête avec le cycle central : ils \textbf{partagent aussi un sommet
entre eux} (et donc une arête entière). On parle de \textbf{fusion coadjacente} ou de structure
\textit{indénoïde}.

\begin{center}
\begin{tabular}{p{7cm}p{7cm}}
  \textbf{Sans coadjacence} & \textbf{Avec coadjacence} \\
  Les deux voisins sont séparés par au moins une arête vide. Chaque voisin introduit ses propres
  sommets externes sans chevauchement. &
  Les deux voisins sont contigus. Le dernier sommet créé par le voisin en position $i$ (adjacent à
  $C_i$) est \textit{réutilisé} comme premier sommet externe du voisin en position $i+1$. \\
\end{tabular}
\end{center}

Concrètement, dans le code, si le cycle en position $\text{pos}$ a un prédécesseur non nul en
position $\text{pos}-1$, alors le sommet d'index 2 du cycle précédent (le premier sommet après
les deux sommets partagés avec le cycle central) est réutilisé, sans en créer un nouveau.

\subsection{Cas du wrap-around (cycle 6 complet)}

Pour un hexagone central dont toutes les 6 positions sont occupées, les positions 0 et 5 sont
aussi coadjacentes. Ce cas ne peut pas être traité dans le passage avant, car la position 5 est
placée \textit{après} la position 0. Un second passage est donc nécessaire après placement de
tous les cycles : on \textbf{fusionne} les deux atomes dupliqués à la jonction pos=0/pos=5 en un
seul atome (méthode \texttt{\_merge\_vertices}).

% ============================================================
\section{Règles de valence : structure de Kekulé}
% ============================================================

\subsection{Principe}

Dans un hydrocarbure aromatique, chaque carbone a une valence totale de 4. Les carbones de cycle
sont tous liés à 2 ou 3 autres carbones par des liaisons C--C (d'ordre 1 ou 2). La valence
restante est complétée par un atome d'hydrogène (ordre 1).

La \textbf{structure de Kekulé} est une représentation avec alternance explicite de simples et
doubles liaisons. Elle n'est pas unique (il peut exister plusieurs Kekulé pour une même molécule),
mais doit vérifier les règles suivantes :
\begin{itemize}
  \item chaque carbone porte \textbf{au plus une} double liaison C=C ;
  \item chaque carbone porte \textbf{au plus un} atome d'hydrogène ;
  \item la \textbf{somme des ordres de liaison} autour de chaque carbone est exactement 4.
\end{itemize}

\subsection{Algorithme de placement des doubles liaisons : couplage maximum pondéré}

Le placement des doubles liaisons est formalisé comme un \textbf{problème de couplage maximum
pondéré} (\textit{maximum weight matching}) sur le graphe carbone--carbone de la molécule.

\subsubsection{Pourquoi pas un algorithme glouton ?}

Une première implémentation utilisait un algorithme glouton qui parcourait les arêtes par
priorité et assignait une double liaison à chaque arête dont aucune extrémité n'était déjà
couverte. Cette approche présentait deux défauts :

\begin{itemize}
  \item Elle \textbf{excluait les arêtes de fusion internes} (les deux extrémités dans le
    cycle central, degré $\geq 3$), alors qu'en chimie ces liaisons peuvent parfaitement être
    des doubles dans une structure de Kekulé (ex. la liaison centrale du naphtalène).
  \item Elle ne garantissait pas le \textbf{maximum de doubles liaisons} : pour 13 carbones
    (ex. \texttt{5\_0\_5\_0\_0\_0\_0}, heptagone central), le glouton trouvait 5 doubles au lieu
    de 6, laissant 3 carbones sans double liaison au lieu d'un seul radical.
\end{itemize}

\subsubsection{Modélisation par couplage maximum}

Le problème est résolu par l'\textbf{algorithme de Blossom} (Edmonds, 1965), implémenté dans
NetworkX via \texttt{nx.max\_weight\_matching(G, maxcardinality=True)} :

\begin{enumerate}
  \item On construit un graphe $G = (V, E)$ où $V$ = ensemble des carbones et $E$ = ensemble
    des liaisons C--C (y compris les liaisons de fusion internes).
  \item Chaque arête reçoit un \textbf{poids} selon le cycle auquel elle appartient :
    \begin{itemize}
      \item arête dans un cycle à 7 : poids $w = 10$
      \item arête dans un cycle à 5 : poids $w = 3$
      \item arête dans un cycle à 6 ou arête de fusion : poids $w = 1$
    \end{itemize}
  \item \texttt{maxcardinality=True} garantit de maximiser \textbf{d'abord le nombre de paires}
    (nombre maximal de doubles liaisons), \textbf{puis le poids total} (préférence pour les
    doubles dans les heptagones).
\end{enumerate}

L'algorithme de Blossom fonctionne sur les graphes \textbf{non bipartites}, ce qui est essentiel
car les pentagones et heptagones (cycles impairs) rendent le graphe carbone non bipartite.
Les approches purement bipartites (matrice de biadjacence, déterminant de Rispoli) ne sont
valables que pour les benzenoïdes purs (cycles pairs uniquement).

\begin{remarque}
  Cette modélisation est équivalente au modèle $M_2^K$ de la thèse de Varet (AMU, 2022) :
  une variable booléenne $x_{v,w} \in \{0, 1\}$ par liaison, avec la contrainte
  $\sum_{w \in \mathcal{N}(v)} x_{v,w} = 1$ pour chaque carbone $v$ (exactement une double
  liaison par carbone). La différence est que Varet utilise PyCSP3 (Choco Solver) pour
  \textit{énumérer toutes} les structures de Kekulé, alors que nous utilisons NetworkX pour en
  trouver \textit{une seule optimale} (poids maximal).
\end{remarque}

\subsubsection{Repli glouton}

Si NetworkX n'est pas installé, le code bascule sur un algorithme glouton corrigé qui considère
\textbf{toutes} les arêtes C--C (y compris les arêtes internes) avec priorité aux arêtes des
cycles à 7. Ce repli est sous-optimal mais fonctionnel.

\subsection{Règle du nombre pair / impair de carbones}

\subsubsection{Formule du nombre de carbones}

Le nombre total de carbones se calcule directement à partir de la séquence. Sans coadjacence :
\[
  N_C = n + \sum_{\substack{i=0 \\ s_i \neq 0}}^{n-1} (s_i - 2)
\]
où $n$ est la taille du cycle central et $s_i$ la taille du voisin en position $i$. Chaque
voisin de taille $s_i$ fusionné sur une arête partage 2 carbones avec le cycle central et en
ajoute $(s_i - 2)$ nouveaux. En cas de coadjacence, un sommet supplémentaire est partagé entre
deux voisins consécutifs, ce qui réduit $N_C$ de 1 par fusion coadjacente.

\begin{remarque}
  L'apport de chaque type de voisin : pentagone $\to +3$ (impair), hexagone $\to +4$ (pair),
  heptagone $\to +5$ (impair). Ainsi, la parité de $N_C$ dépend de la taille du cycle central
  et du nombre de voisins à contribution impaire (pentagones et heptagones).
\end{remarque}

\subsubsection{Nombre pair — couche fermée}

Si $N_C$ est pair, la molécule est dite \textbf{à couche fermée} : le couplage maximum de
Blossom trouve un \textbf{couplage parfait} (chaque carbone est couvert par exactement une
double liaison). Il n'y a \textbf{aucun radical}.

C'est le cas idéal : toutes les valences sont satisfaites à 4, et le champ de force MMFF94
planarise la structure de manière fiable.

\subsubsection{Nombre impair — carbone radicalaire}

Si $N_C$ est impair, il est \textbf{mathématiquement impossible} de former des paires complètes.
Le couplage maximum couvre $N_C - 1$ carbones ($\lfloor N_C / 2 \rfloor$ doubles liaisons),
laissant \textbf{exactement un carbone} sans double liaison : c'est le \textbf{carbone
radicalaire} (un électron célibataire).

\begin{attention}
  Il ne doit y avoir qu'\textbf{un seul} carbone radicalaire, jamais deux ou trois. Si le
  solveur laisse plusieurs carbones sans double liaison, c'est un bug (le couplage n'est pas
  maximal). C'est précisément le problème que l'ancien algorithme glouton produisait.
\end{attention}

Le carbone radicalaire a une valence de \textbf{3} (et non 4). Deux cas possibles :

\begin{center}
\begin{tabular}{lccl}
  \toprule
  Configuration & Liaisons C--C & Hydrogènes & Description \\
  \midrule
  Carbone de jonction  & 3 simples  & 0 H & Carbone interne, valence 3 = 3$\times$C--C \\
  Carbone périphérique & 2 simples  & 1 H & Carbone de bord, valence 3 = 2$\times$C--C + 1$\times$H \\
  \bottomrule
\end{tabular}
\end{center}

\subsection{Procédure de validation chimique}
\label{sec:procedure-chimiste}

La procédure suivante, fournie par le chimiste du projet, doit être respectée pour s'assurer
qu'une structure est chimiquement correcte :

\begin{enumerate}
  \item \textbf{Maximiser l'alternance simple/double} : préférer compléter avec le maximum de
    doubles liaisons dans les cycles à 7, car le champ de force ``planarise'' mieux les
    structures avec des doubles dans les heptagones.
  \item \textbf{Nombre impair $\Rightarrow$ un seul radical} : si $N_C$ est impair, s'assurer
    qu'un unique carbone a une valence de 3 (soit 3 C--C, soit 2 C--C + 1 H). Il faut
    éventuellement \textit{retirer} un H d'un carbone qui a deux liaisons simples avec deux
    autres carbones.
  \item \textbf{Nombre pair $\Rightarrow$ alternance complète} : on doit pouvoir alterner
    simples et doubles liaisons sans aucun carbone radicalaire.
  \item \textbf{C--H dans le plan} : les liaisons C--H doivent rester dans le plan de la
    molécule pour les carbones de valence 4 (garanti par construction dans notre code :
    $z = 0$ pour tous les atomes).
  \item \textbf{Optimisation MMFF94} : le champ de force MMFF94 (et non UFF) est le plus adapté
    pour les petites molécules organiques aromatiques. UFF planarise artificiellement toutes les
    structures (constantes de torsion trop fortes). Si MMFF94 donne une déviation au plan moyen
    $> 10°$, la structure est \textbf{rejetée} comme non planaire.
\end{enumerate}

\begin{remarque}
  Le code Python gère les étapes 1 à 4 automatiquement. L'étape 5 (optimisation MMFF94) est
  réalisée en post-traitement dans un logiciel de visualisation moléculaire (Avogadro, par
  exemple). Le critère de rejet par MMFF94 permet de filtrer les topologies qui ne sont pas
  physiquement réalisables en tant que molécules planes, même si leur graphe est bien construit
  (ex. \texttt{5\_5\_5\_0\_0}, 12 C, planaire avec UFF mais fortement déformée avec MMFF94).
\end{remarque}

\subsection{Placement des hydrogènes}

Pour chaque carbone, on calcule la valence utilisée (somme des ordres de liaison C--C). Si la
valence restante est $\geq 1$, on place un atome H dans la \textbf{direction externe} :
\[
  \vec{u}_{\text{ext}} = -\frac{\displaystyle\sum_{\text{voisins C}} (\vec{r}_{\text{voisin}} - \vec{r}_C)}
                               {\left\|\displaystyle\sum_{\text{voisins C}} (\vec{r}_{\text{voisin}} - \vec{r}_C)\right\|}
\]
La position du H est à $d_{CH} = 1.08\ \text{\AA}$ dans cette direction. Comme tous les carbones
sont dans le plan $z = 0$, les hydrogènes sont automatiquement dans ce même plan : \textbf{la
liaison C--H est planaire par construction}.

% ============================================================
\section{Architecture du code}
% ============================================================

\subsection{Vue d'ensemble}

\begin{lstlisting}[language=bash, style={}]
non_benzenoid_generator/
|-- config.py                 # Constantes physiques et parametres globaux
|-- main.py                   # Point d'entree : boucle sur les tailles de cycle
|-- core/
|   |-- topology.py           # Structures de donnees (MolecularGraph, Vertex, Cycle)
|   |-- geometry.py           # Placement geometrique 3D des atomes
|   `-- valence_solver.py     # Attribution doubles liaisons et hydrogens
|-- generators/
|   `-- builder.py            # Orchestrateur : sequence -> fichier CML
`-- utils/
    |-- symmetry.py           # Reduction par symetrie (groupe diedral)
    `-- exporters.py          # Export Chemical Markup Language (CML)
\end{lstlisting}

\subsection{\texttt{config.py} — Paramètres globaux}

Centralise toutes les constantes physiques et les paramètres de génération :

\begin{lstlisting}
BOND_LENGTH_CC_SINGLE = 1.42  # Angstroms
BOND_LENGTH_CC_DOUBLE = 1.34  # Angstroms
BOND_LENGTH_CH = 1.08         # Angstroms
VALID_CYCLE_SIZES = [5, 6, 7]
NEIGHBOR_VALUES = [0, 5, 6, 7]

NEIGHBOR_CONSTRAINTS = {
    5: {'min': 2, 'max': 4, 'n_positions': 5},
    6: {'min': 2, 'max': 6, 'n_positions': 6},
    7: {'min': 2, 'max': 5, 'n_positions': 7},
}
\end{lstlisting}

\subsection{\texttt{core/topology.py} — Graphe moléculaire}

Définit les structures de données fondamentales :

\begin{itemize}
  \item \texttt{Vertex} : un atome, avec son élément (\texttt{'C'} ou \texttt{'H'}),
    ses coordonnées $(x, y, z)$ et la liste de ses liaisons \texttt{(voisin\_id, ordre)}.
  \item \texttt{Cycle} : une liste ordonnée d'IDs de sommets et une taille.
  \item \texttt{MolecularGraph} : dictionnaire d'atomes + liste de cycles. Méthodes :
    \texttt{add\_vertex}, \texttt{add\_bond}, \texttt{get\_carbon\_degree},
    \texttt{get\_valence\_used}, \texttt{count\_carbons}, \texttt{to\_cml\_data}.
\end{itemize}


\subsection{\texttt{utils/symmetry.py} — Réduction par symétrie}

\textbf{\texttt{get\_dihedral\_transforms(seq)}} génère les $2n$ transformées de la séquence :
\begin{lstlisting}
# Rotations
rot = tuple(seq[(i + k) % n] for i in range(n))
# Reflexions
ref = tuple(seq[(k - i) % n] for i in range(n))
\end{lstlisting}

\textbf{\texttt{canonical\_form(seq)}} choisit la forme canonique (zéros à droite, max lexico).

\textbf{\texttt{generate\_canonical\_sequences(...)}} énumère toutes les combinaisons via
\texttt{itertools.product}, filtre par nombre de voisins, exclut le benzenoïde pur, et ne
conserve qu'une représentante par orbite.

\subsection{\texttt{core/geometry.py} — Moteur géométrique}

\textbf{\texttt{place\_central\_cycle(size)}} :
Place le cycle central en $(0, 0)$, sommets sur le cercle de rayon $R_n$, dans le plan $z=0$.
Crée les liaisons C--C du cycle.

\textbf{\texttt{compute\_cycle\_center(v1, v2, n\_sides)}} :
Calcule le centre du nouveau cycle placé sur l'arête $[v_1, v_2]$, à distance $a_n$ du milieu
de l'arête, du côté extérieur.

\textbf{\texttt{place\_cycle\_with\_shared\_vertex(v1, v2, n\_sides, shared\_vertex\_id)}} :
Place un nouveau cycle fusionné sur $[v_1, v_2]$. Si \texttt{shared\_vertex\_id} est fourni,
le sommet en position $k = n-1$ (adjacent à $v_1$, côté coadjacent) est réutilisé plutôt que
recréé. Les liaisons sont ajoutées sans doublon.

\textbf{\texttt{build\_from\_sequence(central\_size, sequence)}} :
Orchestre la construction complète :
\begin{enumerate}
  \item Place le cycle central.
  \item Pour chaque position non nulle, détecte la coadjacence avec le prédécesseur et place
    le cycle voisin.
  \item \textit{Wrap-around} : si les positions 0 et $n-1$ sont toutes deux non nulles, fusionne
    les deux atomes dupliqués à la jonction.
\end{enumerate}

\textbf{\texttt{\_merge\_vertices(keep\_id, remove\_id)}} :
Fusionne deux sommets : redirige toutes les liaisons de \texttt{remove\_id} vers
\texttt{keep\_id}, supprime \texttt{remove\_id} du graphe.

\subsection{\texttt{core/valence\_solver.py} — Solveur de valences}

\textbf{\texttt{solve()}} : point d'entrée. Appelle dans l'ordre :
\begin{enumerate}
  \item \texttt{\_solve\_matching} (si NetworkX disponible) ou \texttt{\_solve\_greedy} (repli) :
    place les doubles liaisons et retourne l'ensemble des carbones couverts.
  \item \texttt{\_place\_hydrogens} : ajoute les H selon la procédure de validation chimique.
\end{enumerate}

\textbf{\texttt{\_solve\_matching}} (méthode principale) :
\begin{lstlisting}
# Construire le graphe pondere NetworkX
G = nx.Graph()
G.add_nodes_from(carbons)
for vid in carbons:
    for other_id, _ in graph.vertices[vid].bonds:
        if other_id in carbon_set and other_id > vid:
            # Poids selon le cycle d'appartenance
            w = 10 if edge in edges_in_7 else (3 if edge in edges_in_5 else 1)
            G.add_edge(vid, other_id, weight=w)

# Couplage maximum pondere (Blossom)
matching = nx.max_weight_matching(G, maxcardinality=True)
\end{lstlisting}

\textbf{\texttt{\_place\_hydrogens(carbons, covered)}} :
\begin{lstlisting}
for vid in carbons:
    if vid in covered:
        # sp2 : 1H si valence < 4
        if 4 - get_valence_used(vid) >= 1:
            add_hydrogen(vid)
    else:
        # Radical : reserve 1 valence pour l'electron celibataire
        if cc_degree >= 3:
            pass    # jonction, 0H
        else:
            n_h = 4 - cc_degree - 1  # degre 2 -> 1H
\end{lstlisting}

\subsection{\texttt{generators/builder.py} — Orchestrateur}

Assemble le pipeline complet pour une séquence donnée :
\begin{enumerate}
  \item Crée un \texttt{MolecularGraph} vide et un \texttt{GeometryEngine}.
  \item Appelle \texttt{build\_from\_sequence} (topologie + géométrie).
  \item Appelle \texttt{ValenceSolver.solve()} (doubles liaisons + H).
  \item Appelle \texttt{CMLExporter.export()} (fichier \texttt{.cml}).
\end{enumerate}

\subsection{\texttt{utils/exporters.py} — Export CML}

Sérialise le graphe en \textbf{Chemical Markup Language} (XML), format lisible par Avogadro,
OpenBabel, et interprétable par Choco Solver. Chaque atome est exporté avec ses coordonnées 3D,
chaque liaison avec son ordre (1 ou 2).

\subsection{\texttt{main.py} — Point d'entrée}

\begin{lstlisting}
for size in [5, 6, 7]:
    configs = generate_canonical_sequences(
        size, n_positions, max_neighbors, min_neighbors
    )
    for central_size, seq in configs:
        builder.build(central_size, seq, output_dir)
\end{lstlisting}

Les fichiers CML sont écrits dans \texttt{output/cycle5/}, \texttt{output/cycle6/},
\texttt{output/cycle7/}, nommés par leur séquence (ex. \texttt{5\_6\_5\_5\_0.cml}).

% ============================================================
\section{Bibliothèques utilisées}
% ============================================================

Le projet utilise la bibliothèque standard Python~3 et une dépendance externe optionnelle
(NetworkX). Si NetworkX n'est pas installé, le solveur de valences bascule sur un algorithme
glouton de secours.

\begin{center}
\begin{tabular}{lp{9cm}}
  \toprule
  Module & Utilisation \\
  \midrule
  \texttt{math} & \texttt{sin}, \texttt{cos}, \texttt{atan2}, \texttt{sqrt} pour la géométrie \\
  \texttt{itertools.product} & énumération de toutes les combinaisons de séquences \\
  \texttt{xml.etree.ElementTree} & construction de l'arbre XML pour CML \\
  \texttt{xml.dom.minidom} & formatage lisible (indentation) du XML \\
  \texttt{dataclasses} & définition des structures \texttt{Vertex} et \texttt{Cycle} \\
  \texttt{pathlib} & gestion des chemins de fichiers \\
  \texttt{typing} & annotations de types (\texttt{List}, \texttt{Tuple}, \texttt{Optional}, etc.) \\
  \midrule
  \texttt{networkx} & \texttt{max\_weight\_matching} — algorithme de Blossom (Edmonds, 1965)
    pour le couplage maximum pondéré. Installation : \texttt{pip install networkx}. \\
  \bottomrule
\end{tabular}
\end{center}

% ============================================================
\section{État d'avancement}
% ============================================================

\subsection{Fonctionnalités complètes et validées}

\begin{itemize}[label=$\checkmark$]
  \item Énumération canonique sous le groupe diédral $D_n$ pour $n \in \{5, 6, 7\}$
  \item Placement géométrique du cycle central et de tous les cycles voisins
  \item Gestion correcte de la fusion coadjacente (sommets partagés entre voisins consécutifs)
  \item Gestion du wrap-around pour le cycle 6 complet (fusion des atomes dupliqués)
  \item \textbf{Couplage maximum pondéré} (algorithme de Blossom / NetworkX) pour le placement
    des doubles liaisons — maximise le nombre de doubles et préfère les cycles à 7
  \item \textbf{Gestion correcte du carbone radicalaire} : exactement 1 radical si $N_C$ impair,
    0 si pair. Le radical de degré 2 reçoit 1 H (et non 2), le radical de degré 3 reçoit 0 H
  \item Placement des hydrogènes dans le plan de la molécule ($z = 0$ par construction)
  \item Export CML valide pour toutes les structures testées
\end{itemize}

\subsection{Fonctionnalités prévues (travail futur)}

\begin{todo}
\begin{itemize}
  \item \textbf{Optimisation MMFF94} : les structures générées doivent être optimisées avec le
    champ de force MMFF94 (et non UFF qui planarise artificiellement). Si la déviation au plan
    moyen dépasse 10°, la structure est rejetée. Actuellement réalisé manuellement dans Avogadro.
  \item \textbf{Validation de planéité automatique} : calcul du RMSD hors-plan après optimisation
    MMFF94 pour automatiser le filtrage des structures non planaires.
  \item \textbf{Intégration Choco Solver} : utilisation des CML générés comme entrée du modèle
    CSP pour construire la table de voisinage.
  \item \textbf{Détection des isomorphismes} : vérification que deux séquences canoniques
    différentes ne donnent pas accidentellement des graphes moléculaires identiques.
\end{itemize}
\end{todo}

% ============================================================
\section{Exemple complet : la molécule \texttt{5\_6\_5\_5\_0}}
% ============================================================

\subsection{Description}

Cycle central : \textbf{pentagone} (5 carbones). Voisins :
\begin{itemize}
  \item position 0 : pentagone (5)
  \item position 1 : hexagone (6)
  \item position 2 : pentagone (5)
  \item position 3 : pentagone (5)
  \item position 4 : aucun (0)
\end{itemize}

Les positions 0-1, 1-2 et 2-3 sont consécutives et non nulles : \textbf{trois fusions
coadjacentes} sont présentes.

\subsection{Comptage attendu}

\begin{center}
\begin{tabular}{lc}
  \toprule
  Grandeur & Valeur \\
  \midrule
  Atomes C & 19 \\
  Atomes H & 8 \\
  Doubles liaisons C=C & 8 \\
  \bottomrule
\end{tabular}
\end{center}

\subsection{Vérification de la valence}

\[
  \text{Valence totale} = 19 \times 4 = 76
\]
\[
  \text{Liaisons C-C simples} + 2 \times \text{doubles} + \text{liaisons C-H}
\]

Avec 8 doubles, le nombre de liaisons C=C contribue $8 \times 2 = 16$ au comptage de valence,
les simples et les C-H complètent le total de 76. Ce bilan est vérifié par le solveur.

% ============================================================
\section{Glossaire}
% ============================================================

\begin{description}
  \item[HAP] Hydrocarbure Aromatique Polycyclique. Famille de molécules planes composées de
    cycles carbonés fusionnés.
  \item[Non-benzenoïde] HAP contenant au moins un cycle de taille différente de 6.
  \item[Fusion par arête (\textit{edge-fused})] Deux cycles partagent exactement une arête
    (deux carbones et la liaison entre eux).
  \item[Coadjacence / fusion indénoïde] Deux cycles voisins d'un même cycle central qui sont
    aussi directement fusionnés entre eux par une arête.
  \item[Forme canonique] Représentant unique d'une classe d'équivalence sous l'action d'un
    groupe de symétrie.
  \item[Groupe diédral $D_n$] Groupe des symétries d'un polygone régulier à $n$ côtés,
    contenant $n$ rotations et $n$ réflexions.
  \item[Structure de Kekulé] Représentation d'une molécule aromatique avec alternance
    explicite de liaisons simples et doubles.
  \item[Carbone radicalaire] Carbone avec un électron non apparié, résultant d'un nombre
    impair de carbones dans la structure (impossible de former des paires complètes de
    doubles liaisons).
  \item[Algorithme de Blossom] Algorithme d'Edmonds (1965) pour trouver un couplage maximum
    dans un graphe général (non bipartite). Utilisé ici via NetworkX pour placer les doubles
    liaisons de manière optimale.
  \item[Couplage maximum pondéré] Ensemble maximal d'arêtes sans sommet commun, qui maximise
    d'abord le nombre d'arêtes puis la somme des poids. Modélise le placement des doubles
    liaisons dans les structures de Kekulé.
  \item[Couche fermée] Molécule dont tous les électrons sont appariés (nombre pair de carbones,
    couplage parfait). Par opposition au \textit{radical} (nombre impair, un électron
    célibataire).
  \item[UFF] Universal Force Field. Planarise artificiellement les structures — non fiable
    pour valider la géométrie des non-benzenoïdes.
  \item[MMFF94] Merck Molecular Force Field 94. Champ de force recommandé pour les petites
    molécules organiques aromatiques. Critère de rejet : déviation au plan moyen $> 10°$.
  \item[CML] Chemical Markup Language. Format XML standard pour l'échange de données
    moléculaires.
  \item[CSP] Constraint Satisfaction Problem. Problème de satisfaction de contraintes, ici
    résolu avec Choco Solver (Java).
\end{description}

\end{document}
